\bookchapter{Capitolo 64: Archiviazione dei dati e richieste di rete}

\chapterintro{Un'app che serve solo a “guardare” è morta; un'app che salva dati e si collega alla rete è viva. In questo capitolo conosci il “trio” del livello dati di iOS: UserDefaults (le piccole configurazioni), SwiftData (i dati strutturati), URLSession (il dialogo con i server), più Codable per analizzare il JSON e il moderno async/await. A fine capitolo potrai scrivere app vere, capaci di ricordare le abitudini dell'utente e di aggiornare i contenuti in tempo reale.}

\booksection{Come scegliere lo schema di archiviazione: guarda prima questa tabella}

\begin{center}
\begin{tabularx}{\linewidth}{|X|X|X|X|}
\hline
\textbf{Soluzione} & \textbf{Cosa conviene salvare} & \textbf{Caratteristiche} & \textbf{Dove lo usiamo in questo libro} \tabularnewline
\hline
UserDefaults & Configurazioni sparse (interruttori, preferenze) & Coppie chiave-valore, semplice e diretto & Impostazioni del capitolo 65, città preferite del capitolo 66 \tabularnewline
SwiftData & Dati strutturati (voci, registrazioni) & Wrapper ufficiale del database, dichiarato con una riga @Model & Voci di spesa del capitolo 65 \tabularnewline
URLSession & Dati del server (non salvati, presi dalla rete) & Framework di rete di sistema, nessuna dipendenza di terze parti & Dati meteo del capitolo 66 \tabularnewline
\hline
\end{tabularx}
\end{center}

\begin{center}\small\color{gray}La divisione dei compiti del trio dei dati\end{center}

Impara questa filastrocca: “le piccole configurazioni vanno in UserDefaults, i dati in blocco in SwiftData, i dati dall'esterno passano per URLSession”. Core Data dei vecchi tutorial funziona ancora (SwiftData è il suo moderno travestimento), mentre SQLite scritto a mano ormai non si tocca quasi mai: impara direttamente il nuovo, non prendere la strada vecchia.

\booksection{UserDefaults: salvare piccole configurazioni}

\begin{lstlisting}[language=Swift]
// Scrittura: salva un valore
UserDefaults.standard.set(true, forKey: "isDarkMode")
UserDefaults.standard.set("Pechino", forKey: "favoriteCity")

// Lettura: preleva (se assente restituisce il predefinito)
let isDark = UserDefaults.standard.bool(forKey: "isDarkMode")
let city = UserDefaults.standard.string(forKey: "favoriteCity") ?? "Pechino"

// Cancellazione: rimuove una chiave
UserDefaults.standard.removeObject(forKey: "isDarkMode")
\end{lstlisting}

Tieni presente che è adatto solo a piccole configurazioni sparse: decine o centinaia di registrazioni vanno affidate a SwiftData, e non usare UserDefaults come database (rallenterebbe l'avvio e rende scomode le ricerche).

\booksection{SwiftData: dare ai dati una “casa vera”}

SwiftData è la soluzione di database ufficiale dall'iOS 17 e semplifica all'estremo “creare le tabelle e aggiungere, eliminare, modificare, cercare”: con una sola annotazione @Model la classe dei dati diventa una tabella del database:

\begin{lstlisting}[language=Swift]
import SwiftData

// 1. @Model: questa classe diventa automaticamente una "tabella"
@Model
class Record {
    var amount: Double
    var category: String
    var note: String
    var date: Date

    init(amount: Double, category: String, note: String, date: Date) {
        self.amount = amount
        self.category = category
        self.note = note
        self.date = date
    }
}

// 2. Collega il contenitore all'ingresso dell'app
@main
struct MyApp: App {
    var body: some Scene {
        WindowGroup {
            ContentView()
        }
        .modelContainer(for: Record.self)   // database pronto
    }
}

// 3. Nelle view leggi con @Query e scrivi con @Environment
struct RecordListView: View {
    @Query(sort: \Record.date, order: .reverse) var records: [Record]
    @Environment(\.modelContext) private var context

    var body: some View {
        List(records) { r in
            Text("\(r.category) \(r.amount) yuan")
        }
    }

    func add(amount: Double, category: String) {
        let r = Record(amount: amount, category: category, note: "", date: Date())
        context.insert(r)     // inserisce
        try? context.save()   // salva su disco
    }
}
\end{lstlisting}

Questo modello “@Model + modelContainer + @Query” puoi portarlo identico in qualsiasi tuo progetto: basta cambiare i nomi dei campi. Appena i dati cambiano, @Query aggiorna l'interfaccia da solo: una meraviglia.

\booksection{Codable: fare da “interprete” al JSON}

I dati dei server sono per lo più testo JSON. Il protocollo Codable di Swift fa avvenire da sola la conversione “JSON struct”: se i nomi dei campi corrispondono, la magia funziona:

\begin{lstlisting}[language=Swift]
// JSON: {"name":"Pechino","main":{"temp":26.5},"weather":[{"description":"Sereno"}]}
struct WeatherResponse: Codable {
    let name: String
    let main: Main
    let weather: [Weather]
}
struct Main: Codable { let temp: Double }
struct Weather: Codable { let description: String }

// Stringa JSON  struct
let decoder = JSONDecoder()
let data = jsonString.data(using: .utf8)!
let result = try decoder.decode(WeatherResponse.self, from: data)
print(result.weather[0].description)   // "Sereno"
\end{lstlisting}

Regola ferrea: guarda prima il JSON realmente restituito, poi definisci i campi. Molti bug di “analisi fallita” nascono da nomi di campo sbagliati o tipi indovinati a caso (il server manda la stringa “26.5” e tu dichiari un Double): stampa il JSON grezzo, confrontalo e in dieci secondi trovi il colpevole.

\booksection{URLSession + async/await: parlare con il server}

Su iOS la rete si gestisce con URLSession, incluso nel sistema, insieme ad async/await di Swift (le funzioni asincrone): “aspettare la rete” diventa naturale come “aspettare un pacco”:

\begin{lstlisting}[language=Swift]
struct WeatherService {
    // Funzioni asincrone: async significa "devo aspettare un attimo", throws significa "può fallire"
    func fetchWeather(city: String) async throws -> WeatherResponse {
        let url = URL(string:
            "https://api.example.com/weather?city=\(city)")!
        let (data, _) = try await URLSession.shared.data(from: url)
        return try JSONDecoder().decode(WeatherResponse.self, from: data)
    }
}

// Punto di chiamata (dentro .task):
.task {
    do {
        let weather = try await service.fetchWeather(city: "Pechino")
        // aggiorna l'interfaccia...
    } catch {
        print("Qualcosa è andato storto: \(error)")
    }
}
\end{lstlisting}

\textbf{Attenzione alle trappole: le due regole ferree della rete}\par

(1) Le connessioni di rete devono usare async/await e non bloccano mai il thread principale: un'interfaccia bloccata è un'esperienza da crash. La concorrenza strutturata di Swift rende tutto naturale, ma non dimenticare try await e do/catch. (2) ATS vieta per impostazione predefinita l'HTTP in chiaro: iOS esige HTTPS. Per testare un servizio http:// locale devi aprire temporaneamente un'eccezione nell'Info.plist (NSAppTransportSecurity), ma il prodotto ufficiale deve viaggiare tutto in HTTPS: anche la revisione lo controlla.

\begin{tcbpractice}{Mettiti alla prova: crea una piccola app "frase del giorno"}
\booksection{Anche gli errori di rete vanno “classificati”: non scaricare ogni errore addosso all'utente}

Rete assente, timeout, 404, errore 500 del server… Per ogni errore diverso, l'utente dovrebbe leggere messaggi diversi:

\begin{lstlisting}[language=Swift]
func friendlyMessage(_ error: Error) -> String {
    if let urlError = error as? URLError {
        switch urlError.code {
        case .notConnectedToInternet: return "Rete assente — controlli la connessione?"
        case .timedOut:               return "Richiesta scaduta — riprovare più tardi?"
        case .cannotFindHost:         return "Server non trovato"
        default:                      return "Problema di rete: \(urlError.code)"
        }
    }
    if let d = error as? DecodingError {
        return "Formato dati inatteso (l'API potrebbe essere cambiata)"
    }
    return "Qualcosa è andato storto: \(error.localizedDescription)"
}
\end{lstlisting}

Mostra all'utente “parole comprensibili” e registra nei log la “frase originale” (stampa l'error completo). Non scaricare in faccia all'utente il gergo inglese di error.localizedDescription: è in questi dettagli che si nasconde la qualità di prodotto.

Trova un'API gratuita di citazioni (per esempio hitokoto o simili), usa URLSession + async/await per scaricare la “frase del giorno”, mostrala al centro dello schermo e aggiungi un pulsante “Cambia frase”. Poi salva le frasi che ti piacciono nei “preferiti” con SwiftData. Questo esercizio usa insieme UserDefaults, SwiftData, Codable e URLSession: finito questo, i capitoli 7 e 8 saranno un gioco da ragazzi.

\end{tcbpractice}

\begin{tcbsummary}{Riepilogo del capitolo}
\begin{itemize}
  \item Filastrocca: le piccole configurazioni in UserDefaults, i dati in blocco in SwiftData, i dati esterni via URLSession
  \item I tre colpi di SwiftData: @Model dichiara la tabella, modelContainer aggancia il contenitore, @Query legge e context.insert scrive
  \item Codable "traduce" il JSON in automatico: se i campi corrispondono, la magia funziona
  \item URLSession + async/await: stare in rete è come aspettare un pacco; do/catch gestisce i fallimenti
  \item Due regole ferree: non bloccare mai il thread principale; il prodotto ufficiale deve usare HTTPS (ATS)
\end{itemize}

\end{tcbsummary}
